% -----------------------------------------------------------------
% Declaração sobre o uso de ferramentas de inteligência artificial
% -----------------------------------------------------------------
\pretextualchapter{Declaração sobre o uso de ferramentas de inteligência artificial}

Esta dissertação utilizou ferramentas de inteligência artificial generativa em etapas específicas de apoio à elaboração do texto, das figuras e da revisão sistemática da literatura, sempre sob supervisão, conferência e responsabilidade final do autor. Nenhum trecho de prosa final, figura, dado de tabela ou classificação de trabalho correlato foi incorporado de forma inteiramente automática, e todos os elementos foram revisados e validados pelo autor antes da inclusão neste documento.

As ferramentas utilizadas e seus papéis foram os seguintes:

\begin{itemize}
    \item \textbf{Claude (Anthropic), modelo Claude Opus 4.7}: assistência na redação e na revisão gramatical e de concordância, organização da estrutura textual, padronização de convenções tipográficas, geração inicial de tabelas a partir de dados fornecidos pelo autor e conferência dos dados inseridos nas tabelas.

    \item \textbf{ChatGPT (OpenAI), modelo GPT-5.5 thinking high}: segunda camada de verificação na revisão sistemática da literatura, com revisão independente de trabalhos previamente avaliados pelo autor para reduzir vieses de classificação e dar suporte à decisão de inclusão ou exclusão.

    \item \textbf{ChatGPT (OpenAI), modelo de geração de imagens GPT Image 1}: composição parcial de elementos de figuras conceituais. Nenhuma figura foi produzida integralmente por modelo de inteligência artificial. Todas as figuras finais foram compostas, editadas e validadas pelo autor.
\end{itemize}

A responsabilidade pela correção técnica do conteúdo, pela validade das referências bibliográficas, pela coerência entre proposta, método e considerações e pela aderência às normas de redação acadêmica é integralmente do autor.

\cleardoublepage
% ---
